\subsection{The actual characteristic pair of the selected-cusp source loop}
\begin{statement}{P3K.42}{Proved}
Retain the actual positive selected-cusp source loop of P3F.54 and
P3F.56, with fixed sufficiently small scalar, the III-star cap, no
logarithmic modification on the coalescing finite pair, and the
initial physical-delta-edge fan. Its parameter is one positive
\(\rho\)-loop, where \(a(u)=\rho^2\), with the literal double
source closure and the actual Cox normal product. In Jin's endpoint
and physical nonzero-vector-line conventions its characteristic pair and
first-resolution relative coordinates are
\[
       (k_M,\ell_M)=(-6,0),\qquad(k_U,\ell_U)=(6,0),\qquad
       (m_{\rm rel},n_{\rm rel},s)=(0,-6,-1).
\]
The framed local surgery has pair \((5,-4)\). The actual capped
normal sphere bundle in its source-family bordism has pair
\((1,4)\). Their oriented sum gives the meridian-one U pair; the actual physical line M is its Hermitian dual and gives the displayed negative pair.
This statement concerns only the specified initial fan. The complete physical line and Hopf-transfer correction is P3K.68.
\end{statement}
\begin{proof}
\par

\medskip\noindent\textbf{SP.1. The framed local surgery and its actual four-cycle}

\par\nobreak

Let \(e:S^3\hookrightarrow S^6\) be the actual lifted link of P3F.55.
Its projection is the immersion of complex Smale length two, and its
third normal field is the actual increasing height direction. Fix the
first normal field in the oriented complex normal plane of the
projection, and complete the frame \((\nu _1,\nu _2,\nu _3)\) with
that height field. Normalizing and applying Gram--Schmidt uses a
positive triangular homotopy. It therefore retains the actual
projection-normal plus height framing, rather than adding an element
of \(\pi _3(SO(3))\).

Takase, Corollary 2.15, supplies a Seifert embedding
\(V\hookrightarrow S^6\) whose outward boundary normal is this
specified \(\nu _1\). The proof of his Theorem 3.2 gives, for this
projection and this boundary normal,
\[
              2\omega=3\sigma(V)+e(\nu V)^2.             \tag{SP.2}
\]
The invariant and generator comparison in P3F.55, equations
(CL.16)--(CL.19), fixes the convention in (SP.2): the punctured
Kummer test has signature \(-16\) and Smale invariant \(-24\)
in both conventions. Hence the actual local projection has
\(\omega=-2\). The same comparison gives
\[
             8H(e)=e(\nu V)^2-\sigma(V),\qquad H(e)=1.
                                                               \tag{SP.3}
\]
Solving these two displayed equalities gives
\[
                    e(\nu V)^2=5,\qquad \sigma(V)=-3.    \tag{SP.4}
\]

We now identify these numbers with the characteristic numbers of the
\emph{framed} surgery manifold, including its meridian line. Let
\[
 N_e=(S^6\setminus\operatorname{int}(S^3\times D^3))
                  \cup(D^4\times S^2)
\]
use precisely the preceding normal frame. Along the Seifert collar,
the vector from its boundary into \(V\) is \(-\nu _1\).
Consequently the truncated Seifert surface has boundary
\(S^3\times\{-\nu _1\}\), and caps by the literal constant section
\(D^4\times\{-\nu _1\}\). Denote this closed four-manifold by
\(\widehat V\).

Here are the orientations of this cap. Write the ambient orientation
near the knot as
\(\Omega_{S^3}\wedge dn_1\wedge dn_2\wedge dn_3\).
The Seifert orientation is
\(dn_1\wedge\Omega_{S^3}\), so its oriented normal two-plane is
\(-dn_2\wedge dn_3\). At \(-\nu _1\), the outward-first
orientation of the normal meridian sphere is also
\(-dn_2\wedge dn_3\). The boundary of the knot complement is
\(+\Omega_{S^3}\wedge\Omega_{S^2}\); the surgery cap therefore
has orientation \(-\Omega_{D^4}\wedge\Omega_{S^2}\).
The cap of \(\widehat V\) has orientation \(-\Omega_{D^4}\).
Thus its oriented normal plane is exactly the positive meridian
plane, with no sign change in its Euler class.

Let \(x=x_U\) be the meridian class with integral one on this sphere.
It is the Thom-line convention; the physical nonzero-vector line is its
dual, as proved explicitly in SP.2 and SP.5 below.
The Mayer--Vietoris map in degree two identifies the boundary
\(H_2(S^3\times S^2)=\mathbb Z\) with the meridian summands in
both pieces; it gives \(H^2(N_e)=\mathbb Zx\). In degree four,
the connecting map identifies \([\widehat V]\) with the oriented
boundary \([S^3]\). The normal-plane calculation above therefore
gives the actual duality identity
\[
       \operatorname{PD}_{N_e}(x)=[\widehat V],\qquad
       x|_{\widehat V}=e(\nu\widehat V).                 \tag{SP.5}
\]
The cap frame is constant. Its Euler square is the relative square
in (SP.4), and capping by a four-ball leaves the signature unchanged.
Restriction of the ambient tangent bundle to this embedded cycle is
\(T\widehat V\oplus\nu\widehat V\). It follows directly that
\[
 \begin{aligned}
 k(N_e)&=\int_{N_e}x^3=\int_{\widehat V}e(\nu\widehat V)^2=5,\\
 \ell(N_e)&=\int_{N_e}p_1(TN_e)x
       =3\sigma(\widehat V)+e(\nu\widehat V)^2=-4.
 \end{aligned}                                             \tag{SP.6}
\]
This constructs the four-cycle and its Euler-line restriction; it
does not recover the pair from the Haefliger number alone.

\par

\medskip\noindent\textbf{SP.2. The actual Cartier-to-Cox coordinate map}

\par\nobreak

Use the normalization and branch notation of F56:
\[
 N=\mathbf P^1_s\times D_\rho,\quad
 x=s-\rho,\quad y=s+\rho,\quad Z=x/y,\quad
 D_+=(x=0),\quad D_-=(y=0).
\]
Remove the inner four-ball to obtain \(W\). The actual determinant
normal isomorphism of F56 is
\[
 L=\mathcal O(-D_+-D_-),\qquad
 E=\mathcal O(D_+)\oplus\mathcal O(-D_-),\qquad
 \operatorname{ad}E=\mathbb R_\gamma\oplus L_{\mathbb R}.
                                                               \tag{SP.7}
\]
The determinant trivialization is multiplication by \(x/y\).
In the finite \((s,\rho)\) ball, the Cartier frames of \(E\) are
\((1/x,y)\); their determinant is \(y/x\), so the stated
trivialization sends their wedge to one. These are nonvanishing
frames of the indicated line bundles even on the corresponding
divisors.

The morphism \(r=[x:y]:W\to\mathbf P^1\) identifies \(E\)
with \(r^*(\mathcal O(1)\oplus\mathcal O(-1))\). On the chart
\(y\ne0\), the coefficient arrow from the Cartier frames to the
standard dual-tautological and tautological frames is exactly
\[
                         \operatorname{diag}(1/y,y).      \tag{SP.8}
\]
Indeed the first pulled-back divisor frame is \(y/x\), and the
second divisor is absent from this chart. Equation (SP.8) follows by
expressing \(1/x\) and \(y\) in those frames. The corresponding
equation on \(x\ne0\) agrees on the overlap.

The actual positive-weight spinor coefficients are
\((\overline{c_0},c_1)\), while the right-complex Cox coefficients are
\((c_0,\overline{c_1})\). Put \(Y=y/|y|\). The full coefficient map
from (SP.8) is therefore
\[
 (\overline{c_0},c_1)=\operatorname{diag}(\overline Y,Y)(a,b),
 \qquad (c_0,\overline{c_1})
      =\operatorname{diag}(Y,\overline Y)(\overline a,\overline b).
\]
Multiplication by the normalized quaternion \(1+jZ\) gives the actual
anti-linear arrow
\[
 (a,b)\longmapsto g(x,y)\overline{(a,b)},\qquad
 g(x,y)=\frac1{\sqrt{|x|^2+|y|^2}}
 \begin{pmatrix}y&-\overline x\\x&\overline y\end{pmatrix}.
                                                        \tag{SP.9}
\]
It intertwines the positive scalar action on the spinor with the
physical negative right scalar on the Cox quaternion. The previous
linear summand swap omitted this conjugation.

For \(J=\left(\begin{smallmatrix}0&-1\\1&0\end{smallmatrix}\right)\)
and every SU(2) transition, \(\overline g=JgJ^{-1}\). Thus the
global anti-linear map \(Q=J^{-1}C\),
\[
 Q(a,b)=(\overline b,-\overline a),\qquad
 Q(e^{i\theta}v)=e^{-i\theta}Q(v),                       \tag{SP.10}
\]
commutes with all clutches. It has real four-dimensional determinant
one, while its adjoint is the antipodal map of the real normal
three-plane, of determinant minus one. It extends over both cap
normal discs by minus the identity. The Hermitian dual associated-line
map is retained on both cap collars; every extending scalar gauge is
inverted pointwise. The conjugated complex clutch has the same second
Chern number, since conjugation by J is an actual constant bundle map.

The physical meridian sign follows in the literal F9 chart
\((x,v,z)\), with \(y=vz\), weights \((0,-1,+1)\), and
\(b=(|z|^2-|v|^2)/2\). Its positive spinor is
\((\overline v,z)\). At \(v=1,z=0\), the projective coordinate
\(z/\overline v\) has the complex orientation of \(y\), while
the outward raw normal sphere at \(b<0\) has its negative. The
physical circle-first normal orientation is the negative raw
orientation. Thus the physical spinor projective sphere has that
physical normal orientation. Its physical nonzero-vector line is
tautological, with Chern number minus one. The meridian-one Thom
line U is its Hermitian dual; on the source page it is M inverse.
This is a coefficient and circle map, not merely a degree comparison.

The inner projection-normal frame is the Cartier frame used here,
up to an extending gauge. To verify that assertion rather than
postulate it, use F56's explicit nonzero normal covector
\[
       (3s^2-\rho^2)\,dX-2s\,dy
                          +2\rho(s^2-\rho^2)\,d\rho .     \tag{SP.11}
\]
The determinant exact sequence identifies its normal quotient with
the canonical line of the normalization. The ratio of this line
identification and the Cartier identification is a holomorphic unit
on the punctured finite ball. Hartogs extension extends it to the
ball, and a holomorphic function in two variables cannot have its
only zero at the origin. Its extension is therefore a unit. The
resulting real two-plane frame comparison extends over the inner
four-ball. Adding the height line and the explicit triangular
graph-normal shear gives the actual frame of SP.1, with no free
\(SO(3)\) clutch.

The normal-coordinate unit can also be checked directly in both
actual fixed-edge charts. We use the physical source coordinate
\(q\) of F56, not its Tate multiple. With \(X=xy\), restriction
of the differential of that displayed source equation to the
normalization gives
\[
 dq=-\frac{xy}{xy+u(\rho)}\,n,
 \qquad n\text{ is the covector in (SP.11)}.              \tag{SP.11a}
\]
The first Cox chart has \(q=Zc_0c_1\), and the second has
\(q=Z^{-1}c'_0c'_1\), with \(c'_i=Zc_i\). Their quotient
complex normal coordinates consequently satisfy, on the fixed curve,
\[
 \begin{aligned}
 d(c_0c_1)&=-\frac{y^2}{xy+u(\rho)}\,n,\\
 d(c'_0c'_1)&=-\frac{x^2}{xy+u(\rho)}\,n.
 \end{aligned}                                             \tag{SP.11b}
\]
These are exactly the two Hom arrows of (SP.8), multiplied by the
same explicit unit \(-1/(xy+u(\rho))\). Since \(u(0)=u_*\ne0\),
it has an analytic logarithm and square root on the finite inner
ball. Its spinor adjustment therefore extends over that ball. This
verifies the actual normal-unit comparison, as well as its line
transition. It supplies no additional relative Chern number.

There is also an exact check on the separate normal factors. Once
their product has been matched by (SP.11b), any remaining
weight-preserving linear normal gauge has the form
\((c_0,c_1)\mapsto(g c_0,g^{-1}c_1)\). Remove its positive
modulus by the actual equivariant positive rescaling. For its unitary
phase, both right-complex spinor coordinates are multiplied by
\(g/|g|\); hence
\[
                  q_H\longmapsto q_H\,g/|g|.             \tag{SP.11c}
\]
This is the physical circle action, and its adjoint Hopf normal
coordinate is exactly unchanged. It contributes no additional
\(SO(3)\) clutch. A possible parameter winding of this circle-line
gauge is retained and extends over \(D^3\times S^1\), whose
restriction map on \(H^1\) is an isomorphism. We therefore do not
assume that each separate normal gauge has a logarithm across the
parameter origin.

\par

\medskip\noindent\textbf{SP.3. The outer cap: the actual product arrow}

\par\nobreak

The outer coordinate change is the literal map of F56
\[
 A_\rho=\begin{pmatrix}1&-\rho\\1&\rho\end{pmatrix},
 \qquad
 \mathcal B_\rho(P)=
       \frac{A_\rho P A_\rho^*}{\operatorname{tr}(A_\rho P A_\rho^*)}.
                                                               \tag{SP.12}
\]
Here the closed three-ball is the positive semidefinite Hermitian
matrices of trace one. Its boundary consists of rank-one matrices.
This is a periodic diffeomorphism of the ball, including its centre,
and its boundary map is exactly \(r\). We use its actual product
extension on the surgery collar, rather than replacing the normal
frame by its derivative.

The collar in P3K.40 and F56 is the literal first-resolution product
\((h,q_H)\), where \(h\) is the base sphere coordinate and
\(q_H=(1+jZ)(c_0+j\overline{c_1})\) is the normalized actual
Cox quaternion. The physical circle acts by
\(q_H\mapsto q_H e^{-i\theta}\). Thus on the normal sphere
cover the exact outer cap map is
\[
 (P,q_H,\rho)\longmapsto
                  (\mathcal B_\rho(P),q_H,\rho)
 \quad\text{on }D^3\times S^3\times S^1.                 \tag{SP.13}
\]
It is a periodic diffeomorphism, is equivariant for that same circle,
and restricts to the actual base-coordinate change on the boundary.
Its quotient is
\((\mathcal B_\rho,\operatorname{Id}_{S^2})\). It therefore
extends both the actual Cox normal frame and the meridian circle
line over the product cap. This is precisely the product surgery
collar used in F56: its second factor is the actual Hopf normal
sphere, and (SP.13) holds that factor, with its orientation, fixed.

In particular (SP.13) keeps the full determinant winding. For
clarity the exact determinant discrepancy can be written down.
The Hom arrow induced by (SP.8) takes the Cartier frame \(xy\)
of \(L\) to \(y^2\) times the standard \(\mathcal O(-2)\)
frame. Under the standard canonical identification this is
\[
                  y^2\,dZ=2\rho\,ds.                    \tag{SP.14}
\]
Thus the actual Cox arrow has the factor
\(\det A_\rho=2\rho\) relative to the cotangent inverse.
Its winding is one. We neither remove it nor identify (SP.13) with
the tangent derivative of \(\mathcal B_\rho\). It remains in
the Cartier-to-Cox coordinate matrix (SP.9) and hence in the clutch
computed below. The earlier attempted derivative comparison that
omitted this factor is preserved as a superseded unpromoted attempt.

All coordinates in (SP.13) are the actual coordinates already
matched in F56 and P3K.40. In particular we make no additional
replacement of \(q_H\) by a Tate-scaled quaternion. The winding
\(-6\) of \(q_T/q\) is retained by those matched charts.
Equations (SP.11a)--(SP.11b) use the physical \(q\) in those
charts; the nonzero factor \(q_T/q\) remains in the period
identification of the adjacent cell. Replacing that period by one
with factor one would change the marking and is not done here.
The original source and line unit gauges used in F56 extend over
the full source disc. Their displayed logarithms give extending
isotopies on the same normal collar. The induced isotopies of its
boundary extend over the product cap by extending their invariant
vector fields with a collar cutoff; averaging preserves the
physical circle action. This carries these actual gauges in
(SP.13), rather than changing the cap framing.

\par

\medskip\noindent\textbf{SP.4. The relative second Chern number and the normal sphere bundle}

\par\nobreak

Cap \(W\) by the inner four-ball in its projection-normal frame
and the outer \(D^3\times S^1\) in its actual Cox frame. The
resulting base and spinor bundle will be denoted by
\(\overline W\) and \(\overline E\). The base orientation is
\[
 \overline W=(S^2\times D^2)\cup(D^3\times S^1)
                      =\partial(D^3\times D^2).
                                                               \tag{SP.15}
\]
The first piece has its complex orientation, and the second has
minus its displayed product orientation. Thus the boundary of the
outer cap is the negative of
\(K=\mathbf P^1\times S^1\), with \(K\) oriented as the
boundary of \(N\).

Use the conjugate spinor on the inner ball, with its conjugate Cartier
frame, and the actual right-complex Cox frame on the outer cap.
The coefficient map between these is exactly the complex-linear matrix
g in (SP.9), after retaining the global anti-linear map (SP.10).
Its Chern number equals that of E by the explicit J conjugacy.
The resulting bundle extends across the inner ball and is a trivial
\(SU(2)\) bundle over \(N\simeq S^2\). Choose a global frame
agreeing there with the Cartier frame. Such agreement can be made
on the ball because the difference of two extending frames is
itself an extending gauge. Compare this frame with the Cox frame on
\(W\). On the inner boundary its complex-linear coordinate matrix after
that conjugation is exactly g in (SP.9).
Since this comparison extends over \(W\), its degree on the outer
boundary \(K\) is the same as its degree on the inner sphere in
its positive complex-domain orientation.

We calculate that degree with its precise trace convention. The
first column of (SP.9) is \((y,x)\); the map
\((s,\rho)\mapsto(y,x)\) is complex linear with nonzero
determinant \(-2\). Hence its normalized boundary map has degree
one for the complex first-column orientation. At the identity the
positive tangent matrices in that orientation are
\[
 I=\begin{pmatrix}i&0\\0&-i\end{pmatrix},\quad
 J=\begin{pmatrix}0&-1\\1&0\end{pmatrix},\quad
 K_0=\begin{pmatrix}0&i\\i&0\end{pmatrix}.
                                                               \tag{SP.16}
\]
They satisfy \(IJ=-K_0\) and \(\operatorname{tr}(IJK_0)=2\).
Consequently the right Maurer--Cartan form has
\[
           \int_K\operatorname{tr}(dg\,g^{-1})^3=24\pi^2.
                                                               \tag{SP.17}
\]
This explicitly fixes the sign; the quaternion convention with the
opposite group orientation cannot be substituted in (SP.17).

Equation (SP.9) is a coordinate arrow from the frame on \(N\) to
the cap frame. A connection that is zero in the first frame is
therefore \(-dg\,g^{-1}\) in the latter frame near their boundary.
Writing \(\omega=dg\,g^{-1}\), we have
\(d\omega=\omega^2\) and
\[
 \operatorname{tr}\bigl((-\omega)d(-\omega)
                    +\tfrac23(-\omega)^3\bigr)
                    =\tfrac13\operatorname{tr}(\omega^3).
\]
Extend this connection into the cap with a collar cutoff. Stokes'
formula and its negative boundary orientation give the actual
clutch number
\[
 \begin{aligned}
 c_2(\overline E)[\overline W]
    &=\frac1{8\pi^2}\int_{\text{cap}}\operatorname{tr}(F^2)\\
    &=-\frac1{24\pi^2}
                       \int_K\operatorname{tr}(dg\,g^{-1})^3=-1.
 \end{aligned}                                             \tag{SP.18}
\]
For the actual adjoint bundle,
\(p_1(\operatorname{ad}\overline E)=-4c_2(\overline E)\).
This identity follows directly on the diagonal circle: the spinor
has Chern roots \(a,-a\), and the nontrivial real adjoint plane has
Euler class \(2a\). Therefore
\[
             p_1(\operatorname{ad}\overline E)[S^4]=4.    \tag{SP.19}
\]

Let \(B=S(\operatorname{ad}\overline E)\to S^4\), with the
actual positive Hopf meridian class \(x\) from F56. Its determinant
class is zero, so the oriented vertical tangent Euler class is
\(e_v=2x\). The exact splitting
\(T_vB\oplus\mathbb R\cong\pi^*(\operatorname{ad}\overline E)\)
gives \(4x^2=\pi^*p_1(\operatorname{ad}\overline E)\).
Since \(\pi_*x=1\) and \(p_1(TS^4)=0\), we obtain
\[
                         (k(B),\ell(B))=(1,4).             \tag{SP.20}
\]
Its index is zero, as already used in F56. The present calculation
retains the individual numbers of this actual normal bundle.

\par

\medskip\noindent\textbf{SP.5. Both line conventions and the actual physical return}

\par\nobreak

The base orientation is the complex orientation on W. In the actual
F9 normal coordinates the ambient orbit orientation is
\(-\Omega_{\rm source}\wedge d\gamma\), and the physical normal
orientation is \(-d\gamma\wedge\Omega_L\). The anti-linear map
(SP.10) supplies the actual normal reflection between the coefficient
frames. Its extensions over the inner D4 and outer D3 times S1 caps
carry the dual circle character and every retained boundary gauge.
In particular the meridian-one line U extends through both product
surgery collars with exactly the line used in (SP.5) and (SP.20).

The outward boundary of the seven-dimensional quotient bordism is
\(-P_f\sqcup-N_e\sqcup-B\): the first sign is the five-dimensional
orbit followed by the outward parameter radial coordinate, and each
inner sign is the outward boundary of a removed normal disc. Thus
Stokes for U gives
\[
 (k_U,\ell_U)(P_f)=-(5,-4)-(1,4).
\]
The physical family is Jin's mapping torus for the inverse return;
reversing that parameter gives
\[
                (k_U,\ell_U)_{\rm Jin}=(6,0).
                                                        \tag{SP.21}
\]
The raw local complex-plus-height six-dimensional orientation and the
physical orbit orientation are opposite. The outward-meridian-one
Thom line changes sign with that normal orientation; hence the raw
Seifert values (SP.6) enter the U calculation with the signs displayed.
The physical nonzero-vector line M instead satisfies M=U inverse,
so in the same oriented Jin manifold
\[
           (k_M,\ell_M)=-(k_U,\ell_U)=(-6,0),
           \qquad A_M=-1.                              \tag{SP.22}
\]

For completeness the Hopf transfer uses this physical line. The
actual complex exterior is D3\_y times S3\_q, and the sphere is
minus the boundary of D3\_y times D4\_q. Physical M is positive on
the oppositely oriented Hopf base; reverse the y factor at the same
time, so the quotient orientation is unchanged. The F22 positive
Hopf factor is then u=+q and y'=-y. If L=D4\_y times S3\_q and
R=S3\_y times D4\_q, its first and second pieces are -R and -L.
The first-to-second clutch r maps to F(r); rewriting with L first
sets x=F(r), so the second point is F inverse of x. Consequently
lambda\_M(F) is minus Sigma\_(F inverse). The removed ball is
D4\_u times D3\_y', whose boundary is the original complex sphere;
the F8 cut-and-reglue maps coincide on its actual free face and are
the identity on its fixed-tube face. This proves that lambda\_M is
the physical clutched sphere. For U the pieces are +R,+L and the
same clutch gives lambda\_U(F)=Sigma\_(F inverse), its orientation
reverse. All circle gauges are carried by the Hermitian dual map.
These exact arrows are written in mathematical notation in P3K.68.

Therefore the physical source class is -1 modulo 28. The exact
relative-coordinate identities in the physical convention are
\(m=k_M-6A_M\), \(n=6A_M\), giving
\[
                    (m,n,s)=(0,-6,-1).                  \tag{SP.23}
\]
The native central generator still has sign plus one: exchanging
the two Hopf pieces inverts the clutch, and the outer orientation
reversal cancels that inverse on a disc-supported central class.
No value on another fan is assigned here.
\end{proof}
\par

\medskip\noindent\textbf{Audit and computation status}

\par\nobreak

The local Seifert formula is read directly from Takase's Corollary
2.15 and the proof of Theorem3.2. The convention is fixed by the
nonzero Kummer comparison in F55. The earlier derivative matrix was
independently multiplied by root, but its proposed identification
with the actual Cox arrow omitted the determinant factor (SP.14)
and has been withdrawn. The actual product-collar map (SP.13)
retains that factor. Root passed the complete revised proof,
including the Seifert cap cycle and full oriented boundary sum.
The independent SP.2--4 audit passed the actual physical-q normal
units, residual physical-circle gauge, product collar, Cartier-to-Cox
arrow and relative Chern sign.

The certificate is \path{code/part3/source_cusp_characteristic_pair.py}.
It records the derivative matrix only as a diagnostic, and checks
the actual Cartier/Cox product, determinant factor,
the fixed real-line sign, the Maurer--Cartan trace, the two Seifert
equations, and the sphere-bundle and bordism arithmetic. It does not
substitute arithmetic for the bundle and cobordism proofs above.

Physical-line correction, 2026-11-17T15:34:51.130110+00:00: original SP.9/10 proof and certificate preserved at records/k42\_pre\_antilinear\_arrow/. The complete cap and inverse-clutch correction is P3K.68.
